%-------------------------
% Resume in Latex
% Author : Sourabh Bajaj

% Website: https://github.com/sb2nov/resume
% License : MIT
%------------------------

\documentclass[letterpaper,10.5pt]{article}
\usepackage{latexsym}
\usepackage[empty]{fullpage}
\usepackage{titlesec}
\usepackage{marvosym}
\usepackage[usenames,dvipsnames]{color}
\usepackage{verbatim}
\usepackage{enumitem}
\usepackage[pdftex, hidelinks, bookmarks=false]{hyperref}
\usepackage{fancyhdr}
\usepackage{ragged2e}  % for justify
\usepackage{fontawesome} % for mobile symbol
\usepackage{multicol} % for multicolumn lists
\usepackage{xcolor}

\pagestyle{fancy}
\fancyhf{} % clear all header and footer fields
\fancyfoot{}
\renewcommand{\headrulewidth}{0pt}
\renewcommand{\footrulewidth}{0pt}

% Adjust margins
\addtolength{\oddsidemargin}{-0.5in}
\addtolength{\evensidemargin}{-0.5in}
\addtolength{\textwidth}{1in}
\addtolength{\topmargin}{-.25in}
\addtolength{\textheight}{1.0in}

\urlstyle{same}

\raggedbottom
\raggedright
\setlength{\tabcolsep}{0in}

%% MODIFY THE vspace here if you want more space
% Sections formatting
\titleformat{\section}{
  \vspace{-7pt}\scshape\raggedright\large
}{}{0em}{}[\color{black}\titlerule \vspace{-6pt}]

%-------------------------
% Custom commands
\newcommand{\resumeItem}[2]{
  \item\small{
    \textbf{#1}{: #2 \vspace{-1pt}} %-3
  }
}

\newcommand{\resumeSubheading}[4]{
  \vspace{-1pt}\item
    \begin{tabular*}{0.97\textwidth}{l@{\extracolsep{\fill}}r}
      \textbf{#1} & #2 \\
      \textit{\small#3} & \textit{\small #4} \\
    \end{tabular*}\vspace{-1pt} %-5
}

\newcommand{\resumeWorkSubheading}[4]{
  \vspace{-1pt}\item
    \begin{tabular*}{0.97\textwidth}{l@{\extracolsep{\fill}}r}
      \textbf{#1} & #2 \\
    \end{tabular*}\vspace{-1pt} %-8
}


\newcommand{\resumeSubItem}[2]{\resumeItem{#1}{#2}\vspace{-1pt}} %-4

\renewcommand{\labelitemii}{$\circ$}

\newcommand{\resumeSubHeadingListStart}{\begin{itemize}[label=,leftmargin=*, parsep=3pt]}
\newcommand{\resumeSubHeadingListEnd}{\end{itemize}}
\newcommand{\resumeItemListStart}{\begin{itemize}[label=\textbullet]}
\newcommand{\resumeItemListEnd}{\end{itemize}\vspace{0pt}}

%-------------------------------------------
%%%%%%  CV STARTS HERE  %%%%%%%%%%%%%%%%%%%%%%%%%%%%


\begin{document}

%----------HEADING-----------------
% \begin{tabular*}{\textwidth}{l@{\extracolsep{\fill}}r}
\begin{table}
\vspace{-10pt}
\centering
\begin{tabular}{c}
\textbf{\Large Chashi Mahiul Islam} \\
\faPhone \enspace +1 (850) 559-8366 \quad \textbullet \quad \faEnvelope \enspace chashi.buet@gmail.com \\ 

% \quad \textbullet \quad \faEnvelope \enspace ci20l@fsu.edu\\

\faLinkedinSquare \enspace \textcolor{blue}{\url{https://linkedin.com/in/chashi-bd}} \quad \textbullet \quad 
\faGithubSquare \enspace \textcolor{blue}{\url{https://github.com/thechashi}} \quad \textbullet \quad 
\faGlobe \enspace \textcolor{blue}{\url{https://chashi.netlify.app/}}
\end{tabular}\vspace{-10pt}
\end{table}

% Seeking summer 2025 internships; specializing in robust, interpretable AI for enhanced vision-language understanding






%-----------EDUCATION-----------------
\section{Education}
  \resumeSubHeadingListStart
  \resumeSubheading
    {Florida State University}{Tallahassee, FL} 
      {Ph.D. Candidate in Computer Science (CS) (Advisor: Dr. Xiuwen Liu)}{December, 2026 (Expected)}
        \begin{itemize}[itemsep=-3pt, topsep=1pt, label=\textbullet]
            \item Reasearch Interest: Robustness of Large Multimodal Models
        \end{itemize}
    \resumeSubheading
    {Florida State University}{Tallahassee, FL} 
      {M.Sc. in Computer Science}{August, 2025}
        \begin{itemize}[itemsep=-3pt, topsep=1pt, label=\textbullet]
            \item CGPA: 3.97/4.00
        \end{itemize}
    \resumeSubheading
    {Bangladesh University of Engineering and Technology}{Dhaka, Bangladesh} 
      {B.Sc. in Computer Science and Engineering}{April, 2019}
  \resumeSubHeadingListEnd






%--------PROGRAMMING SKILLS------------
\section{Technical Skills}
% \vspace{-15pt}
% \begin{multicols}{2}
 \resumeSubHeadingListStart
    % \resumeSubItem{Languages}{C, C++, Java, Python, R, MATLAB}
    % \resumeSubItem{Parallel programming}{POSIX Threads, OpenMP}
    % \resumeSubItem{Data analysis}{SQL, Excel, Scikit-learn, Pandas, Weka, TensorFlow}
    % \resumeSubItem{Others}{Linux scripting, AWS, \LaTeX}
    

    \resumeSubItem{Languages}{Highly proficient-Python, C/C++, Javascript, SQL, HTML/CSS \textcolor{white}{CVPR, Generative AI and NeurIPS}} % TAILOR
 %  \resumeSubItem{Visualization}{Tableau, Matplotlib, Processing, Bokeh}
    % \resumeSubItem{Frameworks/Libraries/Miscellaneous} {PyTorch, HuggingFace, CLIP, Pandas, NumPy, Scikit-learn, Transformers, NLTK, Keras, Angular, React, Spring Boot, Flask, AWS, Google Cloud, Nginx, Gunicorn, Linux, MacOS, Git, Docker, VS Code  }

%     \resumeSubItem{Frameworks/Libraries/Tools}{PyTorch, TensorFlow, Scikit-learn, HuggingFace, Transformers, CLIP, Prompt Engineering, Pandas, NumPy, NLTK, Matplotlib, Tableau, React, Angular, Flask, Spring Boot, AWS, Google Cloud, Docker, Kubernetes, Git, PostgreSQL, MongoDB, Redis, Linux, macOS, VS Code, Distributed Training
% }

\resumeSubItem{Machine Learning Frameworks/Libraries}{PyTorch, Scikit-learn, HuggingFace, Transformers, CLIP, PEFT, Pandas, NumPy, NLTK, Matplotlib, accelerate, Keras, OpenCV, XGBoost, LangChain, AutoGen, CrewAI, LlamaIndex, LangGraph, Vision-Language Models, Embodied AI, Reinforcement Learning}


\resumeSubItem{Tools}{React, Angular, Flask, Spring Boot, AWS, Google Cloud, Docker, Git, VS Code}

\resumeSubItem{Databases/Systems}{PostgreSQL, MongoDB, Redis, Linux, macOS \textcolor{white}{AI reasoning, symbolic and neuro-symbolic AI methods}} % TAILOR


    
 \resumeSubHeadingListEnd
% \end{multicols}





%-----------EXPERIENCE-----------------

\section{Work Experience}
  \resumeSubHeadingListStart

    \resumeWorkSubheading
      {Software Engineer Intern \textnormal{- Meta, Seattle, WA }}{May 2025 -- Aug 2025}
      {}{}
      \resumeItemListStart
        \item Optimized large-scale ranking systems by implementing feature selection techniques and developing scalable data pipelines to integrate new features, improving key metrics in feed ranking models.  
      \resumeItemListEnd
      \vspace{-8pt}

    \resumeWorkSubheading
      {Graduate Research Assistant \textnormal{- Dept. of CS, Florida State University, Tallahassee, FL }}{Aug 2024 -- Ongoing}
      {}{}
      \resumeItemListStart
        \item Conducting research on robustness and explainability of multimodal AI systems, focusing on vision-language models, large language models, and applying machine learning techniques to enhance AI reliability and security.
      \resumeItemListEnd
      \vspace{-8pt}
    \resumeWorkSubheading
      {Graduate Teaching Assistant \textnormal{- Dept. of CS, Florida State University, Tallahassee, FL }}{Jan 2021 -- Jul 2024}
      {}{}
      \resumeItemListStart
        \item Led recitations, developed assessments, and provided student (60-90) support for various computer science courses, including Secure, Parallel \& Distributed Programming, Data Structures, Advanced Python, and OOP with C++.
      \resumeItemListEnd
      \vspace{-8pt}
       \resumeWorkSubheading
          {Software Developer \textnormal{- CoKreates Limited, Dhaka, Bangladesh } }{Aug 2019 -- Dec 2020}
          {CoKreates Limited, Dhaka, Bangladesh}{}
          \resumeItemListStart
 \item Front-end design and development for Admin Panel and Human Resource Module of ERP system. Implemented RESTful APIs for the backend of the HR module and central logging system. Stack: \textbf{Spring Boot}, \textbf{Angular}
          \resumeItemListEnd
  \resumeSubHeadingListEnd




  %-----------Publications-----------------






\section{Selected Publications}
\begin{itemize}[label={},leftmargin=*, itemsep=-0.1em,]

\item \textbullet \quad \textbf{Islam, C.M.}, Jacob Chacko, S., Nishino, M., Liu, X. (2025). \textbf{NeuroShield-ViT: Mechanistic Understandings of Representation Vulnerabilities and Engineering Robust Vision Transformers.} (\textit{ISVC 2025}, Oral). To appear in Lecture Notes in Computer Science (LNCS), Springer.

% \item \textbullet \quad \textbf{Islam, C.M.}, Mamo, O., Jacob Chacko, S., Liu, X., Yu, W. (2025). \textbf{Spatial-ViLT: Enhancing Visual Spatial Reasoning through Multi-Task Learning.} (\textit{ISVC 2025}, Oral). To appear in Lecture Notes in Computer Science (LNCS), Springer.

\item \textbullet \quad \textbf{Islam, C.M.}, Jacob Chacko, S., Horne, P., Liu, X. (2025). \textbf{Vision Embeddings and Their Role in Hallucination Vulnerabilities of Multimodal Large Language Models.} (\textit{ISVC 2025}, Poster). To appear in Lecture Notes in Computer Science (LNCS), Springer.

\item \textbullet \quad \textbf{Islam, C.M.}, Salman, S., Shams, M., Liu, X., Kumar, P. (2024). \textbf{Malicious Path Manipulations via Exploitation of Representation Vulnerabilities of Vision-Language Navigation Systems.} (\textit{IROS 2024}, Oral). In IEEE/RSJ International Conference on Intelligent Robots and Systems.

% \item \textbullet \quad Rahman, M.I., Emdad, F.B., \textbf{Islam, C.M.}, He, Z. (2023). \textbf{Uncovering the Interplay of Demographics and Healthcare Provider Availability on CMS HCC Risk Scores for Disabled Beneficiaries.} (\textit{EHB 2023}). In International Conference on e-Health and Bioengineering (pp. 593-600). Cham: Springer Nature Switzerland.

% \item \textbullet \quad Rakib, G.A., Adnin, R., Bashir, S.A.A., \textbf{Islam, C.M.}, et al. (2022). \textbf{InnerEye: A Tale on Images Filtered using Instagram Filters-How Do We Interact with Them And How Can We Automatically Identify The Extent of Filtering?} (\textit{MobiQuitous 2022}). In International Conference on Mobile and Ubiquitous Systems: Computing, Networking, and Services (pp. 494-514). Cham: Springer Nature Switzerland.

\end{itemize}







  
%--------  Research Experience  ------------
\section{Research Experience}

\resumeSubHeadingListStart

\resumeSubItem{Neural Expansion: A Unified Mechanism for Deep Network Generalization \textnormal{(Manuscript ready to submit)}}
\item \textbullet \quad Proposed Generalization Intervals framework explaining why more than ~90\% of correct predictions occur within directional stability regions defined by Jacobian singular values. (\href{https://github.com/thechashi/neural_expansion}{\textcolor{blue}{codebase}})
\resumeSubItem{The Illusion of Precision: Numerical Instability of Large Language Models}
\item \textbullet \quad Discovered microscopic input perturbations cause 10-million-fold amplification in LLMs, collapsing logit margins and triggering unpredictable token flips. (\href{https://github.com/thechashi/llm_rounding_error_instability}{\textcolor{blue}{codebase}})


% VISION
\resumeSubItem{Robust Vision Transformers (ViTs) - Understanding and Mitigating Representation Vulnerabilities \textnormal{(ISVC 2025)}}

\item \textbullet \quad Examined vulnerabilities in Vision Transformers (ViTs), revealing how small input changes affect layer-wise representations.

\textbullet \quad Proposed NeuroShield-ViT to neutralize vulnerable neurons and enhance adversarial robustness.


\textbullet \quad Improved classification accuracy by 54.85\% against PGD attacks and 71.6\% against IGO attacks on ImageNet1K.

% VISION-LANGUAGE

\resumeSubItem{Adversarial Exploitation in Robot Vision-Language Navigation \textnormal{(IROS 2024)}}
% {
% \begin{itemize}[itemsep=0pt, topsep=0pt, label=\textbullet]
% \item Developed a novel algorithm to adversarially modify images, exploiting the semantic representation gaps in vision-language models, enabling controlled redirection of robot navigation paths through minimal image alterations. Improved security through a detection mechanism sensitive to noise in manipulated images.
% \item Devised algorithms to adversarially modify images, exploiting semantic representation gaps in vision-language models to redirect robot navigation. Enhanced security via detection mechanism for manipulated image noise. 

\item \textbullet \quad Developed a vision-language navigation (VLN) framework for embodied robotic agents, analyzing perception-to-action pipelines within vision-language-action (VLA) architectures.

\textbullet \quad Studied how representation shifts propagate through navigation policy networks, affecting task planning and real-world embodied decision making.

\textbullet \quad Achieved 91.7\% landmark-conditioned path redirection and proposed a robustness-aware detection mechanism improving navigation policy stability by 96.3\%. (\href{https://github.com/programminglove08/RobustnessVisualNav/tree/main}{\textcolor{blue}{codebase}})

% \end{itemize}}
% \newpage
% VISION-LANGUAGE
\resumeSubItem{Enhancing Visual Spatial Reasoning through Multi-Task Learning \textnormal{(ISVC 2025)}}
% {
% \begin{itemize}[itemsep=0pt, topsep=0pt, label=\textbullet]
% \item Developed a vision-language model incorporating spatial features like depth maps, 3D coordinates, and edge maps through multi-task learning, improving spatial understanding and achieving state-of-the-art results on the Visual Spatial Reasoning (VSR) dataset.
% \item Developed SpatialViLT, a vision-language model using spatial features (depth, 3D coordinates, edge maps) via multi-task learning, enhancing spatial understanding and achieving SOTA on Visual Spatial Reasoning (VSR) dataset.

\item \textbullet \quad Developed SpatialViLT, a vision-language model integrating depth maps, 3D coordinates, and edge representations to improve spatial grounding and cross-modal alignment in embodied intelligence systems.

\textbullet \quad Leveraged multi-task learning to enhance 3D scene understanding for downstream robotic navigation and manipulation, improving Visual Spatial Reasoning benchmark performance by 2.1\%.



% \end{itemize}}

% LANGUAGE
% \resumeSubItem{Adversarial Attacks on Aligned Large Language Models  (LLMs) \textnormal{(Under Review)}}
% % {
% % \begin{itemize}[itemsep=0pt, topsep=0pt, label=\textbullet]

% \item \textbullet \quad Developed an adversarial attack for LLMs, using regularized gradients with continuous optimization to generate valid tokens.


% \textbullet \quad Improved attack success rate by ~60\% and reduced the time taken by 3 orders of magnitude, significantly enhancing efficiency.



% \end{itemize}}


% \end{itemize}}

% VISION
% \resumeSubItem{Inherent Vulnerabilities of Vision Transformers
% for Medical Image Classification \textnormal{(Manuscript ready to submit)}}
% \item \textbullet \quad Explored vulnerabilities in vision transformers, showing that minor changes significantly alter representations.


% \textbullet \quad Developed an efficient detection method achieving 73.1\% accuracy, enhancing model reliability.
% \end{itemize}}

% LANGUAGE - TABULAR



% \resumeSubItem{Multimodal Classification with LLMs: Unstructured Text Feature Extraction \textnormal{(Manuscript ready to submit)}}
% % {
% % \begin{itemize}[itemsep=0pt, topsep=0pt, label=\textbullet]
% % \item Developing text data processing pipelines with private large language models to extract meaningful features with in-context reasoning and developing fusion models to improve classification scores by combining unstructured and tabular data in multimodal settings.
% % \item Created text processing pipelines using private LLMs for feature extraction via in-context learning. Developed fusion models combining unstructured and tabular data to enhance multimodal classification scores.
% \item \textbullet \quad Created text processing pipelines with private LLMs for feature extraction through in-context learning.

% \textbullet \quad Built fusion models combining unstructured and tabular data, boosting multimodal classification accuracy from 76\% to 83\%.

% % \end{itemize}}


% VISION
\resumeSubItem{Leveraging Superresolution for Enhanced Analysis of Low-Resolution Content }
% {
% \begin{itemize}[itemsep=0pt, topsep=0pt, label=\textbullet]
\item \textbullet \quad Evaluated SOTA models for identifying low-resolution deep fakes in collected video data.

\textbullet \quad Improved deep fake detection, boosting test accuracy by 19\% on the FaceShifter dataset with Xception Net. (\href{https://github.com/shuvornb/FakeImageDetection}{\textcolor{blue}{codebase}})
% \end{itemize}}

% VISION
% \resumeSubItem{Retrieval of Brain Tumor with Region Based Image Segmentation \textnormal{(Undergraduate thesis)}}
% % {
% % \begin{itemize}[itemsep=0pt, topsep=0pt, label=\textbullet]
% % \item Developed an automated brain tumor image segmentation pipeline, incorporating Anisotropic Filtering for denoising, contrast enhancement, and skull removal, resulting in a remarkable 11.58\% increase in precision and a minor 0.1\% accuracy improvement when compared to the state-of-the-art method for brain tumor segmentation employing adaptive filtering.
% \item \textbullet \quad Engineered an automated brain tumor segmentation pipeline with Anisotropic Filtering for denoising and skull removal.




% \textbullet \quad Achieved an 11.58\% precision increase and 9.1\% accuracy improvement over the SOTA adaptive filtering method.




% \end{itemize}}
\resumeSubHeadingListEnd











%--------  Selected Projects  ------------
\section{Selected Projects}
 \resumeSubHeadingListStart

 \resumeSubItem{Reliability Limits of Deep De Novo Peptide Sequencing (Transformers, Pyteomics) - 2025}
\item \textbullet \quad Found 73.6\% and 95.8\% prediction instability in Casanovo and $\pi$-PrimeNovo under physics-preserving spectral perturbations, revealing fundamental reliability gaps in de novo sequencing models. (\href{https://github.com/thechashi/denovo-reliability-eval/tree/main}{\textcolor{blue}{codebase}})

 % LANGUAGE
    \resumeSubItem{Deciphering Poorly Formatted Text using Private LLM (Transformers, PEFT, BitsAndBytes) - 2023}{
\item \textbullet \quad Processed unstructured text using private LLMs hosted locally for data privacy, enhancing text processing efficiency. 

\textbullet \quad Applied PEFT (LoRA) for efficient finetuning of summarizers and classifiers, boosting performance with prompt engineering. (\href{https://github.com/thechashi/private_llm}{\textcolor{blue}{codebase}})




}



% TIME SERIES
    \resumeSubItem{Multivariate Time Series Classification (Pytorch, Scikit-learn, tsfresh) - 2023}{
    % \newline  Enhance hospital decision-making in patient care by constructing a predictive model using time-series data and radiology reports to forecast intracerebral hemorrhage outcomes. The XGBoost-based model achieved an accuracy of 74.11\% and precision of 75.51\%, demonstrating its capacity to identify high-risk patients. Data collection was sourced from PubMed. 
\item \textbullet \quad Built a predictive XGBoost model to forecast intracerebral hemorrhage using time-series data and PubMed radiology reports.

\textbullet \quad Achieved 74.11\% accuracy and 75.51\% precision, aiding hospital decisions by identifying high-risk patients. (\href{https://github.com/thechashi/medical_multivariate_timeseries}{\textcolor{blue}{codebase}})}
    
% LANGUAGE
    \resumeSubItem{Comparison of Summarization Techniques Using Language Models ( Transformers, NLTK) - 2023}{
\item \textbullet \quad Investigated transformer models for text summarization, comparing BERT, GPT-2, T5, and BART for medical reports.

\textbullet \quad Addressed input size limits and domain-specific vocabulary, improving summarization effectiveness. (\href{https://github.com/thechashi/Data-Science-Graduate-Projects/tree/main}{\textcolor{blue}{codebase}})}

% DATA MINING
%     \resumeSubItem{Web scraping: Higher Education Funding Offers (Selenium, Scikit-learn) - 2023}{
%     % \newline This project is focused on addressing funding challenges for international students in US higher education. Developed a web scraping pipeline using tools such as requests, Selenium, BeautifulSoup, pandas, and re to extract extensive information on professors from CSRank rankings who are looking for new students. 
% \item \textbullet \quad Addressed funding challenges for international students in US by developing a comprehensive web scraping pipeline. (\href{https://github.com/thechashi/Data-Science-Graduate-Projects/tree/main}{\textcolor{blue}{codebase}})

% \textbullet \quad Used requests, Selenium, BeautifulSoup, and pandas to extract data on professors, from CSRank website, seeking new students.

% }

% CLASSIFICATION
%     \resumeSubItem{Cryptic Species Classification Using SVM (Python, SVM) - 2023}{
%     % \newline Developed a model using Support Vector Machines (SVMs) to classify two cryptic species of shellfish, achieving 97\% accuracy. Explained SVM concepts and mathematical formulations, including handling complex data with slack variables and kernel tricks. Highlighted the effectiveness of SVMs in species classification without genetic analysis.
% \item \textbullet \quad Used an SVM model to classify cryptic shellfish species with 97\% accuracy, demonstrating classification without genetic analysis.

% \textbullet \quad Explained SVM concepts like slack variables and kernel tricks to handle complex data and boost performance. (\href{https://github.com/thechashi/Data-Science-Graduate-Projects/tree/main}{\textcolor{blue}{codebase}})}

% CLASSIFICATION
% \resumeSubItem{Data Preprocessing for Kaggle Dataset (Python, Pandas) - 2023}{
% % \newline Preprocessed a Kaggle dataset, employing various data cleaning methods and feature engineering techniques. Implemented strategies like dropping NaN values, one-hot encoding, and median imputation, resulting in an average accuracy of 0.76. Enhanced features from Name and Cabin variables, leading to improved model performance, with SVM accuracy reaching 0.81. Highlighted the importance of feature engineering and effectiveness of ensemble models. 
% \item \textbullet \quad Preprocessed a Kaggle dataset using a variety of data cleaning and feature engineering techniques, including dropping NaN values, one-hot encoding, and median imputation.

% \textbullet \quad Enhanced key features from Name and Cabin variables, increasing SVM accuracy by 7.3\% and improving overall model performance.

% % \newline Preprocessed Kaggle dataset using various cleaning and feature engineering techniques. Applied data transformation methods to improve model performance. Enhanced key features, increasing SVM accuracy by 7.3\%. 
% (\href{https://drive.google.com/file/d/1BhRydG4OmyObAXpJaxcQ0YB-hBxiLf5-/view?usp=sharing}{\textcolor{blue}{project video}})}

% % DATA ANALYSIS
%     \resumeSubItem{NYC Air Quality Data Analysis (Jupyter Notebook, Pandas, Matplotlib) - 2023}{\newline Analyzed NYC air quality data focusing on pollutants like Nitrogen dioxide, Sulfur dioxide, Ozone, and PM2.5. Demonstrated a positive trend in air quality improvement over the past decade, with significant reductions in Nitrogen dioxide and PM2.5 emissions. Employed data visualization techniques to highlight the impact of emission norms and regulations, and emphasized the importance of maintaining good air quality and continuing pollution reduction efforts. (\href{https://github.com/thechashi/Data-Science-Graduate-Projects/tree/main}{\textcolor{blue}{codebase}})}

% % LANGUAGE AND WEBAPP
%     \resumeSubItem{News Website with a Recommendation System (Flask, Gunicorn,  Nginx, Auth0) - 2022}{
    
% \item \textbullet \quad Built a platform displaying stories from the Hacker News API, suggesting recent stories based on user preferences. (\href{https://github.com/thechashi/hackernews}{\textcolor{blue}{codebase}})

% \textbullet \quad Implemented an LSTM-based classifier on story metadata to predict user preferences, improving recommendation accuracy.

%  }

% VISION
    \resumeSubItem{Image Style Transfer with VGG Network (PyTorch) - 2021}{
    
\item \textbullet \quad Implemented image style transfer methodology using convolutional neural networks. (\href{https://github.com/Chashi-Mahiul-Islam/style-transfer}{\textcolor{blue}{codebase}})

\textbullet \quad Used a pretrained VGG 19 network to extract content and style representations, generating a new target image.
}

% DATABASE
    % \resumeSubItem{Self-Paced Ensemble Classifier (Python, Machine Learning) - 2021}{\newline Developed a Self-Paced Ensemble (SPE) classifier to handle highly imbalanced datasets. Utilized under-sampling techniques to harmonize data hardness and improve classification accuracy. 
    
    % Conducted experiments on datasets including forest cover type, credit fraud, and financial payment simulation, achieving superior performance metrics. Provided comprehensive analysis and codebase for effective imbalanced data classification. 
    
    % (\href{https://github.com/chashi-mahiul-islam-bd/dbproject}{\textcolor{blue}{codebase}}, \href{https://youtu.be/IbS7btS5Yc8}{\textcolor{blue}{video presentation}})}

% TIME SERIES
% \resumeSubItem{Sequence to Sequence Stock Time Series Data Prediction Model (PyTorch) - 2021}{

% \item \textbullet \quad Developed an LSTM-based seq2seq model for stock prediction using Alpha Vantage API and PyTorch. (\href{https://github.com/rsc18/tsprediction}{\textcolor{blue}{codebase}})

% \textbullet \quad Added CLI features for stock prediction, data plotting, and model saving, delivering accurate predictions with documentation.


%  }

% Implemented various functionalities including CLI for stock prediction, data plotting, model saving, and integration with custom datasets. Achieved accurate predictions and provided extensive documentation and usage guidelines. 

 \resumeSubHeadingListEnd






%--------  Courses  ------------
\section{Relevant Courses}
 \resumeSubHeadingListStart
 
    % \resumeSubItem{Graduate}{Computer Architecture, Algorithms, Data Mining, Machine Learning, Stochastic Decision Models}
    
    \resumeSubItem{Graduate}{Introduction to Data Science, Data Mining, Advanced Database, Advanced Algorithms, Analytic Methods, Deep Learning and Reinforcement Learning Fundamentals, Data and Computer Communication}
    
    \resumeSubItem{Undergraduate}{Python, C, C++, OOP, Digital Signal Processing, Computer Graphics, Software Development, Discrete Mathematics, Probability and Statistics, Database Management System, Image Processing and Pattern Recognition}
 \resumeSubHeadingListEnd







%% --------AWARDS AND LEADERSHIP------------
\section{Leadership and Awards}

\resumeSubHeadingListStart
\item \small{
% \vspace{-1pt}
\begin{tabular*}{0.97\textwidth}{l@{\extracolsep{\fill}}r}
\textbf{Naaman Franklin Faile Jr. Graduate Fellowship}: Scholarship for Fall 2025 & 2025 \\
\textbf{1st Runner-up}: ACM at FSU Fall Programming Contest & 2024 \\
\textbf{2nd Runner-up}: ACM at FSU Fall Programming Contest & 2023 \\
\textbf{Cultural Secretary}: Bangladesh Student Assiociation-FSU & 2021 \\
\textbf{President}: BUET Film Society & 2018 \\


\end{tabular*}
\vspace{-5pt}}
\resumeItemListEnd


%-------------------------------------------
\end{document}
